%%%%%%%%%%%%%%%%%%%%%%%%%%%%%%%%%%%%%%%%%%%%%%%%%%%%%%%%%%%%%%%%%%%%%%%%%%%%%%%
%                       CARREGA DE LA CLASSE DE DOCUMENT                      %
%                                                                             %
% Les opcions admissibles son:                                                %
%      12pt / 11pt            (cos dels tipus de lletra; no feu servir 10pt)  %
%                                                                             %
% catalan/spanish/english     (llengua principal del treball)                 %
%                                                                             % 
% french/italian/german...    (si necessiteu fer servir alguna altra llengua) %
%                                                                             %
% listoffigures               (El document inclou un Index de figures)        %
% listoftables                (El document inclou un Index de taules)         %
% listofquadres               (El document inclou un Index de quadres)        %
% listofalgorithms            (El document inclou un Index d'algorismes)      %
%                                                                             %
%%%%%%%%%%%%%%%%%%%%%%%%%%%%%%%%%%%%%%%%%%%%%%%%%%%%%%%%%%%%%%%%%%%%%%%%%%%%%%%

\documentclass[11pt,spanish,listoffigures,listoftables]{tfgetsinf}

%%%%%%%%%%%%%%%%%%%%%%%%%%%%%%%%%%%%%%%%%%%%%%%%%%%%%%%%%%%%%%%%%%%%%%%%%%%%%%%
%                     CODIFICACIO DEL FITXER FONT                             %
%                                                                             %
%    windows fa servir normalment 'ansinew'                                   %
%    amb linux es possible que siga 'latin1' o 'latin9'                       %
%    Pero el mes recomanable es fer servir utf8 (unicode 8)                   %
%                                          (si el vostre editor ho permet)    % 
%%%%%%%%%%%%%%%%%%%%%%%%%%%%%%%%%%%%%%%%%%%%%%%%%%%%%%%%%%%%%%%%%%%%%%%%%%%%%%%

\usepackage[utf8]{inputenc} 
%%%%%%%%%%%%%%%%%%%%%%%%%%%%%%%%%%%%%%%%%%%%%%%%%%%%%%%%%%%%%%%%%%%%%
% Para conseguir que la tabla de contenido no salga en rojo
%%%%%%%%%%%%%%%%%%%%%%%%%%%%%%%%%%%%%%%%%%%%%%%%%%%%%%%%%%%%%%%%%%%%%

\hypersetup{ colorlinks=true, linkcolor=black, urlcolor=cyan, }

%%%%%%%%%%%%%%%%%%%%%%%%%%%%%%%%%%%%%%%%%%%%%%%%%%%%%%%%%%%%%%%%%%%%%%%%%%%%%%%
%                        ALTRES PAQUETS I DEFINICIONS                         %
%                                                                             %
% Carregueu aci els paquets que necessiteu i declareu les comandes i entorns  %
%                                          (aquesta seccio pot ser buida)     %
%%%%%%%%%%%%%%%%%%%%%%%%%%%%%%%%%%%%%%%%%%%%%%%%%%%%%%%%%%%%%%%%%%%%%%%%%%%%%%%



%%%%%%%%%%%%%%%%%%%%%%%%%%%%%%%%%%%%%%%%%%%%%%%%%%%%%%%%%%%%%%%%%%%%%%%%%%%%%%%
%                        DADES DEL TREBALL                                    %
%                                                                             %
% titol, alumne, tutor i curs academic                                        %
%%%%%%%%%%%%%%%%%%%%%%%%%%%%%%%%%%%%%%%%%%%%%%%%%%%%%%%%%%%%%%%%%%%%%%%%%%%%%%%

\title{???? ????????? \\
         ?????????????? ? ?????}
\author{Jorge Rodríguez González}
\tutor{Alfons Juan Císcar}
\curs{2025-2026}

%%%%%%%%%%%%%%%%%%%%%%%%%%%%%%%%%%%%%%%%%%%%%%%%%%%%%%%%%%%%%%%%%%%%%%%%%%%%%%%
%                     PARAULES CLAU/PALABRAS CLAVE/KEY WORDS                  %
%                                                                             %
% Independentment de la llengua del treball, s'hi han d'incloure              %
% les paraules clau i el resum en els tres idiomes                            %
%%%%%%%%%%%%%%%%%%%%%%%%%%%%%%%%%%%%%%%%%%%%%%%%%%%%%%%%%%%%%%%%%%%%%%%%%%%%%%%

\keywords{????, ?????????, ????, ?????????????????} % Paraules clau 
         {?????, ???, ???????????????}              % Palabras clave
         {?????, ????? ?????, ?????????????}        % Key words

%%%%%%%%%%%%%%%%%%%%%%%%%%%%%%%%%%%%%%%%%%%%%%%%%%%%%%%%%%%%%%%%%%%%%%%%%%%%%%%
%                              INICI DEL DOCUMENT                             %
%%%%%%%%%%%%%%%%%%%%%%%%%%%%%%%%%%%%%%%%%%%%%%%%%%%%%%%%%%%%%%%%%%%%%%%%%%%%%%%

\begin{document}

%%%%%%%%%%%%%%%%%%%%%%%%%%%%%%%%%%%%%%%%%%%%%%%%%%%%%%%%%%%%%%%%%%%%%%%%%%%%%%%
%              RESUMS DEL TFG EN VALENCIA, CASTELLA I ANGLES                  %
%%%%%%%%%%%%%%%%%%%%%%%%%%%%%%%%%%%%%%%%%%%%%%%%%%%%%%%%%%%%%%%%%%%%%%%%%%%%%%%

\begin{abstract}
????
\end{abstract}
\begin{abstract}[spanish]
????
\end{abstract}
\begin{abstract}[english]
????
\end{abstract}

%%%%%%%%%%%%%%%%%%%%%%%%%%%%%%%%%%%%%%%%%%%%%%%%%%%%%%%%%%%%%%%%%%%%%%%%%%%%%%%
%                              CONTINGUT DEL TREBALL                          %
%%%%%%%%%%%%%%%%%%%%%%%%%%%%%%%%%%%%%%%%%%%%%%%%%%%%%%%%%%%%%%%%%%%%%%%%%%%%%%%

\mainmatter

%%%%%%%%%%%%%%%%%%%%%%%%%%%%%%%%%%%%%%%%%%%%%%%%%%%%%%%%%%%%%%%%%%%%%%%%%%%%%%%
%                                  INTRODUCCIO                                %
%%%%%%%%%%%%%%%%%%%%%%%%%%%%%%%%%%%%%%%%%%%%%%%%%%%%%%%%%%%%%%%%%%%%%%%%%%%%%%%

\chapter{Introducción}

????? ????????????? ????????????? ????????????? ????????????? ?????????????

\section{Motivación}

????? ????????????? ????????????? ????????????? ????????????? ????????????? 

\section{Objetivos}

????? ????????????? ????????????? ????????????? ????????????? ?????????????

\section{Estructura de la memoria}

????? ????????????? ????????????? ????????????? ????????????? ????????????? 

%\section{Notes bibliografiques} %%%%% Opcional

%????? ????????????? ????????????? ????????????? ????????????? ?????????????

%%%%%%%%%%%%%%%%%%%%%%%%%%%%%%%%%%%%%%%%%%%%%%%%%%%%%%%%%%%%%%%%%%%%%%%%%%%%%%%
%                         CAPITOLS (tants com calga)                          %
%%%%%%%%%%%%%%%%%%%%%%%%%%%%%%%%%%%%%%%%%%%%%%%%%%%%%%%%%%%%%%%%%%%%%%%%%%%%%%%

\chapter{Preliminares}

????? ????????????? ????????????? ????????????? ????????????? ?????????????
%==============================================================================

\section{?? ???? ???? ? ?? ??}

????? ????????????? ????????????? ????????????? ????????????? ?????????????

%==============================================================================

\chapter{MAESTRO}

\section{Introducción y contexto}
Uno de los principales factores que limita el rendimiento de los sistemas de AMT basados en el aprendizaje
profundo es la disponibilidad de datos de entrenamiento etiquetados con alta precisión y en cantidad suficiente
para que los modelos sean capaces de generalizar. Históricamente, los datasets disponibles para esta tarea, como \textbf{MAPS}
(\foreign{english}{MIDI Aligned Piano Sounds}), que no supera las 70 horas de audio, 
lo que suponía una barrera significativa para el entrenamiento de arquitecturas modernas.

En este contexto, \textbf{Google Magenta} publicó en 2018 el dataset \textbf{MAESTRO} (\foreign{english}{Midi and Audio Edited for Synchronous Tracks and Organization}),
que supuso un cambio de paradigma en el campo al ofrecer casi 200 horas de grabaciones de piano de alta calidad con alineación precisa entre audio y etiquetas MIDI.
La publicación de MAESTRO demostró inmediatamente su impacto: el modelo \textbf{Onsets and Frames}, entrenado sobre MAESTRO alcanzó un \textbf{F1 de 67,43\%}
en la evaluación sobre grabaciones Disklavier del conjunto MAPS, superando notablemente el \textbf{50,22\%} obtenido por el mismo modelo 
entrenado sobre MAPS únicamente. Desde entonces MAESTRO se ha consolidado como el estándar de facto para la evaluación y comparación de sistemas de transcripción automática de piano.

En este trabajo se utiliza la versión \textbf{V3.0.0} del dataset, publicada en 2020, que constituye una versión más de depurada y actualmente referencia en la comunidad investigadora.

\section{Origen: International Piano-e-Competition}
Los datos que constituyen MAESTRO provienen de grabaciones recogidas durante diez ediciones de la International Piano-e-Competition, una competición de 
piano de nivel internacional celebrada anualmente en la que pianistas virtuosos interpretan sobre pianos de cola \textbf{Yamaha Disklavier}
equipados con un sistema integrado de captura MIDI de alta precisión. Esta competición es la fuente exclusiva de todas las grabaciones del dataset, lo que dota al 
conjunto de datos de una coherencia excepcional en términos de calidad de la interpretación, condiciones acústicas y equipamiento.

\subsection{Yamaha Disklavier}
El elemento técnico que garantiza la calidad de MAESTRO es el piano Yamaha Disklavier. Es un piano de cola acústico convencional al que se le ha añadido un sistema 
electrónico de captura que registra, con alta precisión temporal, el movimiento de cada tecla y pedal durante la interpretación. A diferencia de los sintetizadores
o controladores MIDI, el Disklavier produce sonido acústico real al mismo tiempo que genera los datos MIDI, lo que garantiza que el audio captado por los micrófonos
y el MIDI registrado sean inherentemente sincrónicos.

La alineación resultante entre la señal de audio y las etiquetas MIDI en MAESTRO es de aproximadamente $\pm 3~\mathrm{ms}$, muy inferior al umbral de percepción
temporal humano (estimado en $20~\mathrm{ms}$), lo que hace que las etiquetas sean consideradas alineadas a nivel fonación. Esta precisión es fundamental para entrenar 
modelos de transcripción: cualquier desalineación sistemática entre el onset de una nota en el audio y su etiqueta MIDI correspondiente introduce ruido en los gradientes
durante el entrenamiento y degrada el rendimiento de los modelos.

\section{Adquisición y anotación}
\subsection{Captura del audio}
El audio se graba en calidad CD o superior, con una frecuencia de muestreo de $44\,1$ a $48~\mathrm{kHz}$ y una profundidad de bit de 16 bits 
PCM\footnote{PCM (\foreign{english}{Pulse Code Modulation}) es una técnica de representación digital de señales analógicas basada en el muestreo y la cuantización, ampliamente utilizada en sistemas de audio sin compresión.} estéreo.
Esta calidad es suficiente para representar con fidelidad toda la gama de frecuencias audibles relevantes para el piano (aproximadamente $28~\mathrm{Hz}$ a $4~\mathrm{kHz}$ para los fundamentales,
con armónicos que se extienden hasta 16-$20~\mathrm{kHz}$). Las grabaciones se almacenan sin compresión con pérdida, lo que preserva toda la información espectral original.

\subsection{MIDI y anotación}
El sistema MIDI del Disklavier captura automáticamente los siguientes eventos para cada interpretación:
\begin{itemize}
\item Velocidades de pulsación de teclas (\foreign{english}{key strike velocities}, canal MIDI, nota, velocidad 0-127): codifican la dinámica de cada nota individual.
\item Posición del pedal sustain (Control Change CC64): el pedal más utilizado en el piano, que prolonga el sonido de una tecla.
\item Posición del pedal sostenuto (Control Change CC66): permite mantener seleccionadas únicamente las notas que están sonando en el momento de su activación.
\item Posición del pedal una corda (Control Change CC67): pedal izquierdo, que reduce el volumen y modifica el timbre.
\end{itemize}

La inclusión de los tres pedales hace a MAESTRO uno de los pocos datasets de transcripción de piano con \textbf{anotación de pedal completa}, lo que resulta
especialmente relevante para sistemas que aspiran a una transcripción expresiva o a la síntesis de audio realista a partir de MIDI.

Cada grabación se segmenta automáticamente en piezas musicales individuales y se anota con los siguientes metadatos:
\begin{itemize}
\item  Compositor canónico (nombre normalizado).
\item Título de la obra.
\item Año de la grabación.
\item Partición sugerida (train/validation/test).
\end{itemize}

%==============================================================================

\chapter{MFT-FiLM}

????? ????????????? ????????????? ????????????? ????????????? ?????????????

\section{?? ???? ???? ? ?? ??}

????? ????????????? ????????????? ????????????? ????????????? ?????????????
%==============================================================================

\chapter{Transkun}

????? ????????????? ????????????? ????????????? ????????????? ?????????????

\section{?? ???? ???? ? ?? ??}

????? ????????????? ????????????? ????????????? ????????????? ?????????????

%%%%%%%%%%%%%%%%%%%%%%%%%%%%%%%%%%%%%%%%%%%%%%%%%%%%%%%%%%%%%%%%%%%%%%%%%%%%%%%
%                                 CONCLUSIONS                                 %
%%%%%%%%%%%%%%%%%%%%%%%%%%%%%%%%%%%%%%%%%%%%%%%%%%%%%%%%%%%%%%%%%%%%%%%%%%%%%%%

\chapter{Conclusiones}

????? ????????????? ????????????? ????????????? ????????????? ????????????? 

%%%%%%%%%%%%%%%%%%%%%%%%%%%%%%%%%%%%%%%%%%%%%%%%%%%%%%%%%%%%%%%%%%%%%%%%%%%%%%%
%                                BIBLIOGRAFIA                                 %
%%%%%%%%%%%%%%%%%%%%%%%%%%%%%%%%%%%%%%%%%%%%%%%%%%%%%%%%%%%%%%%%%%%%%%%%%%%%%%%

\begin{thebibliography}{10}

%%%%%%%%%%%%%%%%%%%%%%%%%%%%%%%%%%%%%%%%%%%%%%%%%%%%%%%%%%%%%%%%%%%%%%%%%%%%%%%
% MODEL D'ARTICLE                                                             %
%%%%%%%%%%%%%%%%%%%%%%%%%%%%%%%%%%%%%%%%%%%%%%%%%%%%%%%%%%%%%%%%%%%%%%%%%%%%%%%
\bibitem{light}
   Jennifer~S. Light.
   \newblock When computers were women.
   \newblock \textit{Technology and Culture}, 40:3:455--483, juliol, 1999.

%%%%%%%%%%%%%%%%%%%%%%%%%%%%%%%%%%%%%%%%%%%%%%%%%%%%%%%%%%%%%%%%%%%%%%%%%%%%%%%
% MODEL DE LLIBRE                                                             %
%%%%%%%%%%%%%%%%%%%%%%%%%%%%%%%%%%%%%%%%%%%%%%%%%%%%%%%%%%%%%%%%%%%%%%%%%%%%%%%
\bibitem{ifrah}
   Georges Ifrah.
   \newblock \textit{Historia universal de las cifras}.
   \newblock Espasa Calpe, S.A., Madrid, sisena edició, 2008.

%%%%%%%%%%%%%%%%%%%%%%%%%%%%%%%%%%%%%%%%%%%%%%%%%%%%%%%%%%%%%%%%%%%%%%%%%%%%%%%
% MODEL D'URL                                                                 %
%%%%%%%%%%%%%%%%%%%%%%%%%%%%%%%%%%%%%%%%%%%%%%%%%%%%%%%%%%%%%%%%%%%%%%%%%%%%%%%
\bibitem{WAR}
   Comunicat de premsa del Departament de la Guerra, 
   emés el 16 de febrer de 1946. 
   \newblock Consultat a 
   \url{http://americanhistory.si.edu/comphist/pr1.pdf}.

\end{thebibliography}
\cleardoublepage

%%%%%%%%%%%%%%%%%%%%%%%%%%%%%%%%%%%%%%%%%%%%%%%%%%%%%%%%%%%%%%%%%%%%%%%%%%%%%%%
%                           APÈNDIXS  (Si n'hi ha!)                           %
%%%%%%%%%%%%%%%%%%%%%%%%%%%%%%%%%%%%%%%%%%%%%%%%%%%%%%%%%%%%%%%%%%%%%%%%%%%%%%%

\APPENDIX

%%%%%%%%%%%%%%%%%%%%%%%%%%%%%%%%%%%%%%%%%%%%%%%%%%%%%%%%%%%%%%%%%%%%%%%%%%%%%%%
%                         LA CONFIGURACIO DEL SISTEMA                         %
%%%%%%%%%%%%%%%%%%%%%%%%%%%%%%%%%%%%%%%%%%%%%%%%%%%%%%%%%%%%%%%%%%%%%%%%%%%%%%%

\chapter{Configuració del sistema}

????? ????????????? ????????????? ????????????? ????????????? ?????????????

\section{Fase d'inicialització}

????? ????????????? ????????????? ????????????? ????????????? ?????????????

\section{Identificació de dispositius}

????? ????????????? ????????????? ????????????? ????????????? ?????????????

%%%%%%%%%%%%%%%%%%%%%%%%%%%%%%%%%%%%%%%%%%%%%%%%%%%%%%%%%%%%%%%%%%%%%%%%%%%%%%%
%                               ALTRES  APÈNDIXS                              %
%%%%%%%%%%%%%%%%%%%%%%%%%%%%%%%%%%%%%%%%%%%%%%%%%%%%%%%%%%%%%%%%%%%%%%%%%%%%%%%


\chapter{??? ???????????? ????}

????? ????????????? ????????????? ????????????? ????????????? ????????????? 



%%%%%%%%%%%%%%%%%%%%%%%%%%%%%%%%%%%%%%%%%%%%%%%%%%%%%%%%%%%%%%%%%%%%%%%%%%%%%%%
%                              FI DEL DOCUMENT                                %
%%%%%%%%%%%%%%%%%%%%%%%%%%%%%%%%%%%%%%%%%%%%%%%%%%%%%%%%%%%%%%%%%%%%%%%%%%%%%%%

\end{document}
